\documentclass[11pt]{article}
\usepackage[margin=0.86in]{geometry}
\usepackage[T1]{fontenc}
\usepackage{amsmath,amssymb,amsthm}
\usepackage{graphicx}
\usepackage[hidelinks]{hyperref}
\usepackage{microtype}

\newtheorem{theorem}{Theorem}
\newtheorem{lemma}[theorem]{Lemma}
\newtheorem{corollary}[theorem]{Corollary}
\newtheorem{remark}[theorem]{Remark}
\newcommand{\R}{\mathbb R}
\newcommand{\Var}{\operatorname{Var}}

\title{A Derivative Obstruction to Almost Periodicity in Variation}
\author{fa\_banach\_001 / agent\_lane\_00}
\date{17 August 2026}

\begin{document}
\maketitle

\begin{center}
\fbox{\parbox{0.92\textwidth}{
\textbf{Claimed status: full negative answer to the question in Example 3.3
of arXiv:2205.09031, subject to human review.}
For every $1\leq p<\infty$, the displayed function is Stepanov-$p$ almost
periodic, as recalled in the source, but it is \emph{not}
Stepanov-$p$ almost periodic in variation.
}}
\end{center}

\section{Source question and definitions}

Set
\begin{equation}\label{eq:f}
 f(t)=\sin\!\left(\frac{1}{\zeta(t)}\right),\qquad
 \zeta(t)=2+\cos t+\cos(\sqrt2t).
\end{equation}
The source asks whether $f$ is $S^p$-almost periodic in variation, meaning
that its Bochner transform
\[
 \widehat f(t)\in L^p([0,1]),\qquad \widehat f(t)(s)=f(t+s),
\]
is almost periodic in variation.

\begin{figure}[htbp]
\centering
\includegraphics[width=0.91\textwidth]{figures/open_question_crop.png}
\caption{The question on PDF page 19 of arXiv:2205.09031v1.}
\end{figure}

For a continuous map $H:\R\to Y$ into a Banach space, write
\[
 \Var_1(H;q)=\sup_{q-1=t_0<\cdots<t_N=q+1}
       \sum_{j=0}^{N-1}\|H(t_{j+1})-H(t_j)\|_Y.
\]
The source's definition, specialized to ordinary almost periodicity, says
that $H$ is almost periodic in variation if, for every $\varepsilon>0$,
the set of $\tau$ satisfying
\begin{equation}\label{eq:apvar}
 \sup_{q\in\R}\bigl(\|H(q+\tau)-H(q)\|_Y
  +\Var_1(H(\cdot+\tau)-H(\cdot);q)\bigr)<\varepsilon
\end{equation}
is relatively dense.

\section{The transfer lemma}

\begin{lemma}[Variation transfers to the derivative]\label{lem:transfer}
Let $F\in C^1(\R)$ and $1\leq p<\infty$.  If
\(\widehat F(t)=F(t+\cdot)\) is almost periodic in variation as an
$L^p([0,1])$-valued map, then $F'$ is Stepanov-$1$ almost periodic.
In particular,
\begin{equation}\label{eq:s1bound}
       \sup_{q\in\R}\int_q^{q+1}|F'(u)|\,du<\infty.
\end{equation}
\end{lemma}

\subsection*{Idea of the proof}

The variation of a differentiable Banach-valued curve is the integral of
the norm of its derivative.  Differentiating a translation difference of
the Bochner transform produces translation differences of $F'$.  A short
Tonelli argument then shows that the former variation dominates the usual
Stepanov-$1$ distance between the latter functions.

\begin{proof}
Put $G=F'$ and, for a fixed $\tau$, define
\[
 H_\tau(r)=\widehat F(r+\tau)-\widehat F(r).
\]
On every compact $r$-interval this is a $C^1$ curve in $L^p([0,1])$, with
\[
 H_\tau'(r)(s)=G(r+s+\tau)-G(r+s).
\]
Consequently,
\begin{equation}\label{eq:variationidentity}
 \Var_1(H_\tau;q)=\int_{q-1}^{q+1}\|H_\tau'(r)\|_{L^p([0,1])}\,dr.
\end{equation}
For every $q\in\R$, Tonelli's theorem and
$\|h\|_1\leq\|h\|_p$ on the unit interval give
\begin{align}
 \int_q^{q+1}|G(u+\tau)-G(u)|\,du
 &\leq \int_{q-1}^{q+1}\int_0^1
       |G(r+s+\tau)-G(r+s)|\,ds\,dr \notag\\
 &\leq \int_{q-1}^{q+1}\|H_\tau'(r)\|_p\,dr
  =\Var_1(H_\tau;q).                                      \label{eq:key}
\end{align}
Indeed, after the change $u=r+s$, the double integral has convolution
weight exactly one for $u\in[q,q+1]$.

Thus every $\varepsilon$-translation in \eqref{eq:apvar} is a
Stepanov-$1$ $\varepsilon$-translation of $G$, and such translations are
relatively dense.  This proves Stepanov-$1$ almost periodicity.

For completeness, boundedness follows directly.  Take $\varepsilon=1$ and
let $L$ witness relative density.  Given $x$, choose a $1$-translation
$\tau\in[x-L,x+L]$.  Then
\[
 \int_x^{x+1}|G(u)|\,du
 \leq 1+\int_{x-\tau}^{x-\tau+1}|G(v)|\,dv,
\]
and the final integral is bounded because $x-\tau\in[-L,L]$ and $G$ is
continuous.  This proves \eqref{eq:s1bound}.
\end{proof}

\section{The specific function has unbounded derivative variation}

The denominator in \eqref{eq:f} is strictly positive.  Equality to zero
would force both cosines to equal $-1$, making $\sqrt2$ a ratio of odd
integers.  Hence $f\in C^\infty(\R)$ and
\begin{equation}\label{eq:g}
 g(t):=f'(t)=
 \frac{\sin t+\sqrt2\sin(\sqrt2t)}{\zeta(t)^2}
 \cos\!\left(\frac1{\zeta(t)}\right).
\end{equation}

\begin{lemma}\label{lem:unbounded}
There is a fixed $b\in(0,1)$ and a sequence of intervals
$[m_n,m_n+b]$ such that
\[
             \int_{m_n}^{m_n+b}|g(t)|\,dt\longrightarrow\infty.
\]
In particular, $g$ is not Stepanov-$1$ bounded.
\end{lemma}

\subsection*{Proof intuition}

The irrational orbit $(t,\sqrt2t)$ comes arbitrarily close to
$(\pi,\pi)$ modulo $2\pi$.  Near each such visit, $\zeta$ has a positive,
strictly convex well whose minimum tends to zero.  On one side of the well,
the change of variables $u=1/\zeta(t)$ turns the total variation of
$\sin(1/\zeta)$ into a long integral of $|\cos u|$.

\begin{proof}
The sequence $(2k+1)\sqrt2$ is dense modulo $2$.  We may therefore choose
integers $k_n$ such that, for $x_n=(2k_n+1)\pi$,
\[
 e^{i\sqrt2x_n}\longrightarrow-1.
\]
It follows that
\[
 \zeta(x_n)\to0,\qquad \zeta'(x_n)\to0,\qquad
 \zeta''(x_n)\to3.
\]
More precisely, uniformly for $|h|\leq a$,
\[
 \zeta''(x_n+h)\longrightarrow \cos h+2\cos(\sqrt2h).
\]
Choose a fixed, sufficiently small $a\in(0,1/2)$.  For all large $n$,
\begin{equation}\label{eq:convex}
                  \zeta''(t)\geq1\quad(|t-x_n|\leq a).
\end{equation}
Since $\zeta'(x_n)\to0$, strict convexity gives a unique
$m_n=x_n+o(1)$ in this interval with $\zeta'(m_n)=0$.  It is the local
minimum, and
\[
       0<\delta_n:=\zeta(m_n)\leq\zeta(x_n)\longrightarrow0.
\]
Fix $b=a/2$.  For large $n$, \eqref{eq:convex} holds throughout
$[m_n,m_n+b]$; hence $\zeta$ is increasing there and
\[
              \zeta(m_n+b)\geq\delta_n+\frac{b^2}{2}
              \geq\frac{b^2}{2}.
\]
Using \eqref{eq:g} and then the monotone substitutions
$v=\zeta(t)$ and $u=1/v$, we obtain
\begin{align*}
 \int_{m_n}^{m_n+b}|g(t)|\,dt
 &=\int_{\delta_n}^{\zeta(m_n+b)}
       \frac{|\cos(1/v)|}{v^2}\,dv\\
 &=\int_{1/\zeta(m_n+b)}^{1/\delta_n}|\cos u|\,du
   \longrightarrow\infty.
\end{align*}
The lower endpoint is at most $2/b^2$, whereas the upper endpoint tends
to infinity, and $|\cos u|$ has positive mean $2/\pi$.  Since $b<1$,
these intervals lie inside unit intervals, proving the last assertion.
\end{proof}

\section{Answer}

\begin{theorem}[Full negative resolution]\label{thm:answer}
For every $1\leq p<\infty$, the function in \eqref{eq:f} is not
$S^p$-almost periodic in variation.
\end{theorem}

\begin{proof}
If it were, $\widehat f$ would be almost periodic in variation as an
$L^p([0,1])$-valued map.  Lemma~\ref{lem:transfer} would make $f'=g$
Stepanov-$1$ bounded, contradicting Lemma~\ref{lem:unbounded}.
\end{proof}

\begin{remark}[Scope]
The argument does not dispute the source's statement that $f$ is ordinary
Stepanov-$p$ almost periodic.  It separates that property from the strictly
stronger variation version.  Lemma~\ref{lem:transfer} is also a general
obstruction: any $C^1$ function with unbounded local $L^1$ derivative
variation gives the same negative conclusion for every finite $p$.
\end{remark}

A bounded novelty search of the run registry, exact question wording,
displayed formula, and the 2025 journal publication found no later explicit
answer.  Novelty is provisional; the proof above is self-contained.

\begin{thebibliography}{9}
\bibitem{CKV}
B. Chaouchi, M. Kosti\'c, and D. Velinov,
\emph{Metrical approximations of functions}, arXiv:2205.09031v1 (2022);
Funkcialaj Ekvacioj \textbf{68} (2025), 1--44,
\href{https://doi.org/10.1619/fesi.68.1}{doi:10.1619/fesi.68.1}.
\end{thebibliography}

\end{document}
